% Tema metropolis com paleta própria.

\usetheme[numbering=fraction,progressbar=frametitle]{metropolis}

\usepackage{graphicx}
\usepackage{booktabs}
\usepackage{array}
\usepackage{microtype}
\usepackage{xspace}
\usepackage{caption}
\usepackage{xcolor}

% Paleta: mAccent é a cor de contraste dos diagramas.
\definecolor{mPrimary}{HTML}{0A7E3B}
\definecolor{mPrimaryLight}{HTML}{7FCF88}
\definecolor{mPrimaryDark}{HTML}{064F26}
\definecolor{mNeutral}{HTML}{5B6770}
\definecolor{mAccent}{HTML}{B45309}

\setbeamercolor{structure}{fg=mPrimary}
\setbeamercolor{title separator}{fg=mPrimary}
\setbeamercolor{progress bar}{fg=mPrimary,bg=mPrimary!20}
\setbeamercolor{progress bar in head/foot}{fg=mPrimary,bg=mPrimary!20}
\setbeamercolor{progress bar in section page}{fg=mPrimary,bg=mPrimary!20}
\setbeamercolor{frametitle}{fg=white}
\setbeamercolor{alerted text}{fg=mPrimaryDark}
\setbeamercolor{example text}{fg=mPrimary}
\setbeamercolor{block title}{bg=mPrimary,fg=white}
\setbeamercolor{block body}{bg=white,fg=black}
\setbeamercolor{block title alerted}{bg=mPrimaryLight!45,fg=mPrimaryDark}
\setbeamercolor{block body alerted}{bg=white,fg=black}
\setbeamercolor{itemize item}{fg=mPrimary}
\setbeamercolor{itemize subitem}{fg=mPrimaryDark}
\setbeamercolor{itemize subsubitem}{fg=mNeutral}
\setbeamercolor{enumerate item}{fg=mPrimary}
\setbeamercolor{caption name}{fg=mPrimaryDark}

\metroset{
  sectionpage=progressbar,
  titleformat=smallcaps,
  block=fill,
  numbering=fraction
}
